% Source: https://artofproblemsolving.com/community/c3h2320592_for_once_i_will_not_do_cp_maybe_a_bit
% Stable identity: aops-topic-2320592-q02
\textbf{(a)} Any such expression first combines two of the numbers and then combines the result with the third. The possible results of the first step are:
\begin{itemize}
\item from $3,4$: $7,\ \pm1,\ 12,\ \dfrac34,\ \dfrac43$; combining with $1$ gives at most $12+1=13$ in absolute value;
\item from $3,1$: $4,\ \pm2,\ 3,\ \dfrac13$; combining with $4$ gives at most $4\cdot4=16$;
\item from $4,1$: $5,\ \pm3,\ 4,\ \dfrac14$; combining with $3$ gives at most $5\cdot3=15$.
\end{itemize}
(In the three cases the first-step results have absolute value at most $12$, $4$ and $5$. Combining such a result $x$ with the third number $y$ by $+$, $-$ or $\times$ gives at most $|x|+|y|$ or $|x||y|$ in absolute value, that is at most $13$, $16$ and $15$. Dividing gives at most $12$: $\dfrac xy$ is at most $|x|$, and the only divisions by a number less than $1$ in absolute value are $1\div\dfrac34=\dfrac43$, $4\div\dfrac13=12$ and $3\div\dfrac14=12$.) Every result has absolute value at most $16<24$, so $24$ is impossible.

\textbf{(b)} For example, $\boxed{3\times4\times(1+1)=24}$, or $4\times(3+3)\times1=24$.
